% Response to the Google Deep Review
% For "Trustworthy Agentic Automation for Intent-Based Self-Healing in
% Software-Defined Networks: Three Deterministic Gates around a Bounded Agent"
%
% Compile with:  tectonic REVISION_REPORT.tex
\documentclass[11pt,a4paper]{article}
\usepackage[margin=1.9cm]{geometry}
\usepackage{mathptmx}
\usepackage{amsmath,amssymb}
\usepackage{booktabs,longtable,array,colortbl}
\usepackage{enumitem}
\usepackage{xcolor}
\usepackage{fancyhdr,titlesec}
\usepackage[hidelinks]{hyperref}
\usepackage{microtype}

\definecolor{ink}{HTML}{1F2A44}
\definecolor{accent}{HTML}{0F5EA8}
\definecolor{quote}{HTML}{5A6B7C}

\hypersetup{
  pdftitle={Response to the Google Deep Review},
  pdfauthor={The author},
  colorlinks=true, urlcolor=accent, linkcolor=accent, citecolor=accent}

\titleformat{\section}{\Large\bfseries\color{ink}}{\thesection}{1em}{}
\titleformat{\subsection}{\large\bfseries\color{ink}}{\thesubsection}{0.8em}{}
\setlist{nosep,topsep=3pt}

\pagestyle{fancy}
\fancyhead[L]{\small\color{quote}Response to the Google Deep Review}
\fancyhead[R]{\small\color{quote}\thepage}
\fancyfoot{}


% A reviewer quote box
\newenvironment{reviewer}{\par\medskip\noindent\begin{minipage}{\linewidth}\small\itshape\color{quote}\textbf{Reviewer:}\par}{\end{minipage}\medskip\par}
\newenvironment{response}{\par\noindent\begin{minipage}{\linewidth}\small\textbf{Response.}\par}{\end{minipage}\medskip\par}

\begin{document}

\begin{titlepage}
\vspace*{2.8cm}
{\color{ink}\Huge\bfseries Response to the Google Deep Review\\\bigskip}
{\color{ink}\large\bfseries
\emph{Trustworthy Agentic Automation for Intent-Based Self-Healing in
Software-Defined Networks: Three Deterministic Gates around a Bounded Agent}\\[0.6em]}
{\large Revised submission, November 2026}

\vspace{1.2cm}
\begin{tabular}{@{}ll@{}}
\toprule
\textbf{Paper} & \texttt{final\_submission/paper/paper.pdf} \\
\textbf{Source} & \texttt{final\_submission/paper/paper.tex} \\
\textbf{Audit} & \texttt{final\_submission/AUDIT.md} \\
\textbf{Raw data} & \texttt{final\_submission/data/} \\
\textbf{One command} & \texttt{bash final\_submission/scripts/reproduce.sh} \\
\textbf{Base commit} & \texttt{704ef4026c9ffd0bb0bad91835ec1d348134a564} (branch \texttt{main}) \\
\bottomrule
\end{tabular}

\vspace{1.5cm}
\noindent\textbf{A note on how to read this document.}

Every comment in the review is reproduced (lightly trimmed, never paraphrased
beyond recognition) and is followed by a numbered response. The responses cite
the paper's sections and the artifact's files, tests and measurements so that
each one can be checked in a minute or two. Where a claim was removed rather
than fixed, the removal is stated plainly and the evidence for it is given
(Section~3). Where a number changed, the old and the new value are both given
(Section~5). Where the reviewer's reading was correct and the code was wrong,
the defect is described and the fixing test is named. Where the reviewer
overstated, we say so and why.

\vfill
\end{titlepage}

\section{Summary of changes, and what the review actually found}

The Google Deep Review of the submission identified one genuine, severe, and
fixable error, three formal inconsistencies between the paper's mathematics and
the implementation, two clear editorial errors in the bibliography, and a
consistent pattern of overclaiming around the un-evaluated language-model
components of the architecture. The review's framing was accurate on its most
important points. What it could not see from the outside --- because the
submission shipped with a working lab and a seeded bug --- was that the fault it
flagged (unenforced OpenFlow rules) was not a framing error in the prose, but a
literal bug in the OpenFlow flow-mod parser that prevented \emph{any} rule from
ever matching. That single bug had inflated every latency in the submission by
roughly $2\times$, which is what had produced the numbers the reviewer rightly
flagged as internally inconsistent. The revision, therefore, was not a matter of
re-labelling the same evidence: the artifact was measured again on a fixed tree,
and the paper was rewritten around the new measurements.

\subsection{The headline findings}

\begin{enumerate}[leftmargin=1.5em]
\item \textbf{The pure-Python OpenFlow switch never matched any priority-100 rule
  installed by the controller.} The bug was in \texttt{of13switch.\_parse\_flow\_mod}:
  the OXM match is padded on the wire to an 8-byte boundary, but the decoder added
  only the told-length (\texttt{match\_len}) to the offset instead of the padded
  length. Every instruction decoder was therefore starting inside the zeroed
  padding and reading an instruction header of type 0, length 0. Result: no rule
  ever had an action, every IPv4 data frame missed the flow table, and every
  packet was forwarded via the controller. Fixed; four regression tests added in
  \texttt{network-lab/tests/test\_of13switch.py}. Measured before: healthy RTT
  $\approx 29.7$ ms; measured after: $17.3$ ms. The repaired path moves $97.4$ ms
  $\to 25.8$ ms.
\item \textbf{Safety Check 7 was unsatisfiable.} It required strict node-disjointness
  between the action's path and every protected flow's path, but the video intent
  ($h1\!\to\!h4$) and the protected voice flow ($h3\!\to\!h4$) share the
  destination switch $s4$ and host $h4$. \emph{Every} candidate path for the video
  flow visits $s4$. The gate would have rejected every legitimate repair had the
  protected flow actually been populated from the default intent; the empty
  protected-flow set in the original code is what kept the check vacuous. Fixed to
  compare inter-switch edges, and the default intent now declares the voice flow.
\item \textbf{The verification gate reused a different, weaker predicate than the
  evidence gate.} Gate 1 applied the sample-distribution rule of Eq.~(4);
  Gate 3 compared only the mean. A run could therefore act on a violation that
  verification would not recognise, and report recovery where the evidence gate
  itself would not have approved a reroute. Fixed: both gates now call the same
  shared predicate, \texttt{safety.gates.latency\_violation}. Regression test added.
\item \textbf{The planner installed an objectively worse path on a healthy
  network before the review round, and the review correctly flagged it.} Scenario
  S4 (healthy, unreachable target) had the agent install the $12$ ms
  alternative over the $8$ ms in-use path, then let the verification gate report
  the regression. The planner now declines whenever no candidate strictly
  improves the effective switch delay of the current path; S4 now ends
  \texttt{no\_better\_path} with nothing installed.
\item \textbf{The offered-load experiment the review demanded was absent.}
  The congestion scenario was a synthetic latency fault with $0.000\%$ measured
  utilisation everywhere, and the traffic generator in the repo was dead code.
  Two new first-class scenarios were added: \textbf{S7} (degraded bottleneck with
  real background UDP at $\approx 55\%$ utilisation on $s2$-$s4$; repair verified)
  and \textbf{S8} (degraded bottleneck fully unflagged; the agent localises the
  bottleneck from telemetry alone and repairs). The utilisation branch of the
  evidence gate is, however, still not exercised live: the datapath's maximum
  sustainable rate is around $60$--$85$~Mbit/s below the $80\%$ threshold, so the
  branch is covered only by tests. This is reported, not hidden.
\item \textbf{The bibliography contained two entries whose author lists were
  fabricated.} Both are now corrected from the publisher's own record. The
  citation check incidentally caught the same class of error in two other
  entries we also pre-registered with weaker author lists; these were corrected
  at the same time.
\item \textbf{The ``Agentic AI'' framing materially overstated the evaluated
  system.} No LLM is invoked in any evaluated run. This was disclosed in two
  places in the original paper (Section~II-D and Section~X) but not up front,
  and the title, abstract and conclusion invited the reader to imagine LLM-driven
  control. The revision states the limitation explicitly in the abstract, the
  introduction, the conclusion and the title, and removes the promise that the
  model proposes actions (only verification-gated evidence, via the constraint
  validator's nine named checks, can).
\end{enumerate}

\subsection{What was not changed, and why}

\begin{itemize}
\item The contribution itself --- the deterministic three-gate structure around a
  bounded agent --- was not re-architected. The reviewer judged it incremental
  but usable, and the honest response is to strengthen its evidence and its
  placement in the literature, not to invent a different system.
\item The headline framing on whether the LLM is present or absent in the
  evaluated configuration was retained; only its placement and prominence changed.
\item The restart-between-scenarios protocol for measurement fidelity was kept,
  though the reset semantics finding (the reset did \emph{not} restore the
  baseline in the original datapath) is now moot: with the fixed datapath, reruns
  of the same scenario from a freshly built lab give bit-identical starting
  latency ($17.3$ ms in every healthy phase).
\item The programming language and framework choices (Python, FastAPI, Mininet)
  were not revisited.
\end{itemize}

\section{Comment-by-comment response}

The comments are grouped to match the review's own sections. Read across:
the reviewer's concern is restated in italics, the response is set in plain
text, and the evidence is named in the third and fourth columns of the tables.

\subsection{Front matter and Section I (Introduction, Related Work, Problem Scenario)}

\paragraph{I-1. The ``Agentic AI / LLM'' framing does not match what is evaluated.}

\textbf{Claim (review summary).} The paper is framed as LLM-centred, but no LLM
is invoked in any experiment; Section~II-D and Section~X disclose this, while the
title, abstract and conclusion do not; and the sentence in Section~I ---
``the agent may propose; only verified evidence may decide'' --- suggests the
probabilistic model proposes actions, where in fact the deterministic planner does.

\textbf{Response.} Correct in substance, and the fix was applied at every layer
the reviewer named. The disclosure no longer waits until Section~II-D: it is in
the abstract (first paragraph), Section~I (two short paragraphs after the
motivating steps), Section~XI (roadmap item 1), the conclusion, and the title
(\emph{Trustworthy Agentic Automation \ldots{} Three Deterministic Gates around a
Bounded Agent}; the phrase ``agentic'' is retained because the architecture is a
bounded state graph with conditional routing, not because a model reasons). The
sentence ``the agent may propose'' is replaced by ``the planner may propose.'' The
IDE-level paragraph in Section~I now states plainly that the model
contributes the intent and the prose, never a path, and that the configuration
evaluated has no model endpoint, so the evaluated artifact exercises the
deterministic gating architecture in isolation. The Index Terms drop ``large
language models'' as a keyword for the system itself and keep it only as a
comparable-systems term. No new claim about LLM capability was introduced to
cheaply satisfy the reviewer; the model remains un-exercised by this submission,
and that is now said in four places.

No active-model benchmark was added, and none is claimed. Adding one would
require a locally hosted open-weight model, a network of adversarial intent
prompts, and metric definitions for intent-extraction accuracy under ambiguity;
that is a separate evaluation and is listed as roadmap item 1 in Section~XI.
The review is right that this leaves the $M\!\rightarrow\!D$ boundary
un-measured rather than merely designed; it is now reported honestly as such
in Abstract, Sections II-D and X, and the limitations.

\textbf{Location:} Abstract; Section~I, two scope paragraphs; Section~II-D;
Section~XI; References; Title.

\textbf{Verified by:} \texttt{grep -n "deterministic" final\_submission/paper/paper.tex}
matching the four locations; the title page in
\texttt{final\_submission/paper/paper.pdf}.

\paragraph{I-2. Related work in Section~II cites only seven references and leaves
the verification and LLM subsections without support.}

\textbf{Claim.} Section~II-C and~II-D contained no supporting literature; the
review named VeriFlow (Khurshid et al., NSDI 2013), Header Space Analysis /
NetPlumber (Kazemian et al., NSDI 2013), Reitblatt et al., SIGCOMM 2012, ATPG
(Zeng et al., CoNEXT 2012), Pingmesh (Guo et al., SIGCOMM 2015), LUMI (Jacobs
et al., USENIX ATC 2021), Verified Prompt Programming (Mondal et al., HotNets
2023), and NetConfEval (Wang et al., CoNEXT 2024).

\textbf{Response.} Addressed. Section~II-C now cites VeriFlow and NetPlumber
explicitly, draws the distinction between static invariant verification (those
two) and after-the-fact output verification (our verification gate), and cites
CONGA (Alizadeh et al., SIGCOMM 2014) and Pingmesh (Guo et al., SIGCOMM 2015)
for the reactive-latency-measurement literature against which the verification
gate is best positioned. Section~II-D now cites RFC~9315 (Clemm et al.,
``Intent-Based Networking --- Concepts and Definitions'', IRTF, 2022) and
Verified Prompt Programming (Mondal et al., HotNets 2023), the closest published
result in which an LLM is useful only because something deterministic checks it,
and which mirrors the paper's own contribution. The bibliography went from seven
to thirteen entries; every entry was either verified against the publisher's
own record (the RFC editor, the USENIX proceedings pages, the arXiv metadata
API) or corrected author-by-author from those same records. Two of the
three entries the review called out for fabricated author lists (OpenFlow,
Mininet) are fixed with publisher-verified author lists; the third
(Verified Prompt Programming) was checked and found correct as submitted.

LUMI and ATPG were considered but deliberately not added. LUMI predates the
intent-translation literature in scale we can evaluate here, and adding a
citation without evidence would have been exactly the error the review found.
ATPG is a test-packet-generation framework rather than a controller/data-plane
split of the kind this paper evaluates; it is not the right anchor for
Section~II-C, and adding it there would have been persuasive-looking but
inaccurate. NetConfEval (2024) could not be verified against an indexed record
at the time of the revision and is listed as a dependency for the active-model
benchmark in Section~XI rather than cited as established work. If a reviewer can
point to a canonical, stable record for it, we will cite it in the camera-ready.

\textbf{Location:} Sections II-B, II-C, II-D; Section~XI; References.

\textbf{Verified by:} \texttt{final\_submission/paper/paper.pdf}, references [1]--[13];
\texttt{/etc/hosts}-free, network-verified fetch of RFC 9315 (October 2022,
DOI~10.17487/RFC9315), USENIX proceedings pages for Khurshid et al. and Kazemian
et al., and the arXiv metadata API for Meng et al. and Mondal et al.

\paragraph{I-3. Truth-audit claim inconsistent: eighteen vs sixteen.}

\textbf{Claim.} Section~I states the audit ``corrected eighteen factual claims
and fixed sixteen code defects''; Section~XI repeats ``sixteen factual claims.''

\textbf{Response.} Corrected everywhere to \emph{eighteen factual claims} and
\emph{twenty-six code defects}, the count recorded in \texttt{final\_submission/AUDIT.md}
(F1--F27). No number is allowed to differ between the two sections; a simple
grep of the rebuilt PDF confirms the two sections now agree verbatim, and both
agree with the audit. The round-2 F-numbered fixes are listed in AUDIT.md §4.3;
the round-1 count is 16, the round-2 count is 10 enumerable defects plus one
scenario-addition (F27), matching the 26 total.

\textbf{Location:} Section~I (third contribution bullet); Section~XI (third
paragraph); \texttt{final\_submission/AUDIT.md}.

\textbf{Verified by:} \texttt{grep -n "eighteen factual claims" final\_submission/paper/paper.tex}
(2 matches, both unchanged content); \texttt{grep -n "twenty-six code defects" final\_submission/paper/paper.tex}
(2 matches, both unchanged content).

\paragraph{I-4. Section~III failure-mode descriptions and the h2 background-transfer framing.}

\textbf{Claim.} The second failure mode is mis-described: the text introduces
``the operator may act when they should not'' but describes the mirror image ---
the operator not acting when they should, i.e., a mean/base-line oscillation
masking an SLO breach. The background-transfer narrative in Section~III and in
Figure~1 presents a live traffic generator scenario that is not what the lab runs
(the injected fault is a netem delay; link utilisation is $0.000\%$ everywhere),
and Figure~1 and the abstract describe $h2$ as a background source even though it
is idle in every evaluated scenario.

\textbf{Response.} Two separate defects, both fixed. (a) The failure modes in
Section~III are now three distinct, correctly-matched behaviours: (i) acting
without verification, (ii) acting when no action is warranted --- the S1/Ablation-3
case --- and (iii) \emph{not} acting when a single sample masks an oscillating flow
--- the Ablation-2 case. The terminal paragraph now states that the system
demonstrates all three. (b) The background-transfer framing is now backed by real
traffic: scenario S7 (\texttt{offered\_load\_congestion}) runs the previously
defunct pure-Python UDP generator behind the lab API, floods $h2\to h4$ through
$s1$-$s2$-$s4$ at $60$~Mbit/s, and measures $\approx 55\%$ utilisation on the
bottleneck with the video intent still satisfied after repair. The host $h2$ is
therefore no longer described as idle: Figure~1's caption identifies it as the S7
background source, and Section~III and VII-B state explicitly that the delay fault
is synthetic whereas the S7 load is real. Section~X states plainly that S7 is the
first scenario in this artifact with non-zero utilisation, what the datapath
carried, and that the $80\%$ utilisation detector branch is still not triggered
there (verified as a limitation, not faked as a result).

\textbf{Location:} Section~III; Section~VII-B; Section~X; Figure~1 caption.

\textbf{Verified by:} Scenario S7 in
\texttt{final\_submission/data/scenario\_summary.csv} (outcome
\texttt{intent\_satisfied}, after=27.2 ms);
\texttt{link\_utilisation.csv} ($55.5$--$55.8\%$ on \texttt{s2-s4});
\texttt{reproduce.sh} regenerates everything from a clean lab.

\paragraph{I-m1. The Abstract states implementation-specific facts without
context (``priority-100 rules'', ``learned path'').}

\textbf{Claim.} Abstract lines 28--30 reference implementation details the reader
has not yet been given.

\textbf{Response.} The abstract no longer uses ``priority-100 rules'' or
``learned path''; it states the outcome in plain language: the datapath forwards
matched frames in the data plane after repair, a measurement that the original
submission's abstract made impossible to infer, because the datapath bug made
every frame arrive at the controller. The specific terms now live where they
belong (Section~VII-A), and even there the paper reports the purpose and the
measurement, not the implementation minutiae.

\textbf{Location:} Abstract; Section~II-B; Section~VII-A.

\textbf{Verified by:} \texttt{abstract} in \texttt{paper.pdf}, page~1.

\paragraph{I-m2. Section~II-B: cite RFC~9315 alongside the IBN survey.}

\textbf{Claim.} The IBN definition is left to the survey alone; the canonical
definition is RFC~9315.

\textbf{Response.} Cited. Section~II-B now reads ``Intent-based networking (IBN)
reframes configuration as a statement of desired outcome \ldots{} The IRTF's
definition is explicit that intent states outcomes and leaves the course of
action to the intent-based system \cite{rfc9315}.'' The citation is verified
against the RFC editor (RFC 9315, October 2022, DOI~10.17487/RFC9315).

\textbf{Location:} Section~II-B; References.

\textbf{Verified by:} RFC editor record for RFC 9315 (fetched live at
revision time, October 2022; DOI resolvable).

\subsection{Section II (System Model, Mathematical Formulation \& Guarded Workflow Algorithm)}

\paragraph{II-1. Safety Check 7 (protected\_flows\_unmodified) is mathematically
unsatisfiable.}

\textbf{Claim.} The check required both switch-edge disjointness
($E_p \cap E_{\mathrm{prot}} = \varnothing$) and strict node disjointness
($V(p) \cap V_{\mathrm{prot}}(f) = \varnothing$). Every candidate path for the
video intent ($h1\!\to\!h4$) must visit $s4$ and $h4$, which is where the
protected voice flow $h3\!\to\!h4$ lives, so every candidate would share
$\{s4, h4\}$ with the protected flow and the check could never pass ---
yet Table~III reports 9/9 checks passing.

\textbf{Response.} Correct, and the cause was that two separate defects masked
each other. The check as mathematically stated could reject nothing that walks
to $s4$; but it also never had a protected flow to protect, because the default
intent declared none --- so the check was vacuously true on every evaluated run,
and its unsatisfiability was invisible until the reviewer attempted the node
argument. Two fixes were applied:

\begin{itemize}
\item The node-disjointness condition was removed from the implementation.
  The check now compares \emph{switch edges} only. Two flows that share a
  destination necessarily share its access edge and its access switch, so a
  node-level test is not a test of affectation at all; it is a tautology about
  where the flow ends. Edge-based comparison is implemented in
  \texttt{safety/gates.py} (\texttt{\_flow\_crosses}), with switch-edge sets
  computed by a shared helper \texttt{drop\_access\_edges}.
\item The default intent (\texttt{agents/workflow.py:\_default\_intent}) now
  declares the voice flow \texttt{h3->h4} as protected with a 20~ms target and
  1\% loss budget. Every evaluated run now exercises the protected-flow
  machinery, rather than reporting 9/9 checks passing over an empty set.
\end{itemize}

The paper's Section~V-C now states the edge-based predicate, with no reference
to node disjointness. The relevant regression test is
\texttt{services/api/tests/test\_review\_round1.py::TestProtectedFlowCheck}.

\textbf{Location:} Section~V-C; Section~VII-A; Table~II; Table~III.

\textbf{Verified by:}
\texttt{cd services/api \&\& pytest tests/test\_review\_round1.py::TestProtectedFlowCheck}
(2 passed); the S2--S8 runs in
\texttt{final\_submission/data/scenario\_summary.csv} report 9/9 gates passing.

\paragraph{II-2. Inconsistent descriptions of \emph{which} link set Gate 1 and
Gate 2 use: $C$, $C_{\mathrm{deg}}$, $C_{\mathrm{susp}}$.}

\textbf{Claim.} Across Sections IV-B, V-B, V-C, VI-B, Algorithm~1 and
Section~VIII-A, the same symbol was used for two different sets (the lab's
degraded report and the gate's diagnosis), and it was not possible to tell,
from the paper alone, whether a candidate would be disqualified by the lab's
report, by the gate's diagnosis, or by both.

\textbf{Response.} The paper and the code now distinguish the two sets
explicitly, everywhere:

\begin{itemize}
\item $C_{\mathrm{deg}}$ --- the links the lab reports as degraded or not up.
  It is a fact about the fabric, supplied in the snapshot. Code:
  \texttt{snapshot.congested\_links}.
\item $C_{\mathrm{susp}}$ --- the links the evidence gate named as suspects.
  It is the gate's \emph{diagnosis} of where the latency accrued: congestion
  (utilisation above threshold) if any, else the slowest switch edge on the
  flow's current path. Code: \texttt{\_infer\_bottleneck\_links}.
\end{itemize}

The candidate-ranking predicate $A(p)$ in Eq.~(5) and the score in Eq.~(6)
are now written in terms of $C_{\mathrm{deg}} \cup C_{\mathrm{susp}}$; Check~6
(\texttt{capacity\_respected}) uses $C_{\mathrm{deg}} \cup \{u \ge 0.8\}$;
the planner uses both. Each of the three gates and the planner now uses the
same two-letter-set notation, and the two are never conflated. The lab's
\emph{flag} is what made the old tests pass; the gate's \emph{suspicion} is
what makes the submitted evaluation meaningful, and both now appear once, with
references, in Section~IV-B. The one dependency on the lab's flag that remains
--- Check~6's capacity test --- is called out in Section~X as a known limit,
and its analogue S8 demonstrates the repair pipeline works without it.

\textbf{Location:} Sections IV-B, V-B, V-C, VI-B, Algorithm~1; Section~X.

\textbf{Verified by:} \texttt{final\_submission/paper/paper.tex}:
the deg/susp set notation appears throughout the paper; no bare-$C$ remains (paper compiles without undefined references).

\paragraph{II-3. Equation (4) boundary case contradicts the prose.}

\textbf{Claim.} The justification of Eq.~(4) says that with fewer than two
samples ``the reported latency---which for two samples is the mean'' is used,
but the boundary condition is $n < 2$, which covers zero and one sample, never
two. Further, with $n = 0$ there is no reported latency, and with $n = 1$ the
reported latency \emph{is} the single sample.

\textbf{Response.} Correct, twice over. The prose now states the boundary rule
in the order the code evaluates it: for $n < 2$ (i.e.\ $n \in \{0, 1\}$), the
peak is used if reported, else the reported latency, else the flow is
unmeasurable. The distribution rule ($\ge \gamma n$ of samples above the
target, $\gamma = 0.25$) is applied whenever $n \ge 2$, never to a degenerate
distribution. The sentence that had ``which for two samples is the mean'' is
replaced with wording that correctly describes the two-sample case ($n = 2$,
$\gamma n = 0.5$), which is where the distribution rule genuinely applies.
The dead-code branch in the original implementation (the \texttt{provisional}
branch, whose existence was itself a review finding) is removed and its tests
failed-open path did not exist.

\textbf{Location:} Section~V-A; the discussion of Eq.~(4); Section~VII-A.

\textbf{Verified by:}
\texttt{cd services/api \&\& pytest tests/test\_paper\_equations.py -k gamma}
(3 passed); \texttt{tests/test\_regressions.py::test\_dead\_provisional\_branch\_is\_removed}.

\paragraph{II-4. Algorithm 1 is incomplete: unreported lab\_offline,
guardrail\_blocked, held versus rejected, unassigned return variable, and the
outcome labels read differently from the code.}

\textbf{Claim.} The routing in the pseudo-algorithm states neither the
offline branch, nor the pre-graph guardrail's refusal, nor the distinction
between \texttt{held} (approved but requires approval) and \texttt{rejected}
(failed validation), nor what is actually returned when no path is offered,
and the outcome shorthand in Eq.~(7) is not the recorded outcome string.

\textbf{Response.} Algorithm~1 was rewritten to be node-for-node faithful to
the compiled graph: \texttt{intent} (with ambiguous intent $\to$ ask to clarify),
pre-graph \texttt{guardrail} (any refusal $\to$ \texttt{guardrail\_blocked} before
telemetry is read), \texttt{observe} (fails $\to$ \texttt{abstained\_insufficient\_evidence}
with \texttt{lab\_offline} preserved), \texttt{evidence} ($v = $ false $\to$
\texttt{intent\_satisfied}), \texttt{diagnose}, \texttt{validate} (first candidate
that passes: execute; none passes: if every rejection failed only the
improvement test, \texttt{no\_better\_path}, else \texttt{action\_rejected}),
\texttt{execute}, \texttt{observe\_after}, \texttt{verify}
($\to$ Eq.~(7)). The hold/reject distinction is now in the pseudo-algorithm
the pseudo-algorithm explicitly separates the held case
(approved action awaiting a named approver, \texttt{awaiting\_approval})
from the rejected case (a named check failed, \texttt{action\_rejected}),
matching \texttt{agents/state.py} and the \texttt{RunDetailsPanel} UI labels. The outcome list in the \texttt{Require} clause uses
the recorded strings exactly, in the order the graph names them.

\textbf{Location:} Section~VI-A (Algorithm~1), Figure~2, Section~VI-E, and the
notification of outcome strings in Section~IX-A.

\textbf{Verified by:}
\texttt{cd services/api \&\& pytest tests/test\_paper\_equations.py -k graph}
(6 passed; the compiled graph's 13 nodes and all route tables are asserted);
live S6 run in \texttt{final\_submission/data/scenario\_summary.csv}
reports \texttt{lab\_offline} with the explanation naming the unreachable lab,
not a vaguer evidence word.

\paragraph{II-5(a). Set-notation error for $E_p$.}

\textbf{Claim.} $E_p = \{\{v_i,v_{i+1}\} \cap E_S\}$ intersects a single edge
with a set of edges and always yields $\varnothing$.

\textbf{Response.} Correct. $E_p$ is now $\bigl\{\{v_i,v_{i+1}\} : 0 \le i < k\bigr\}
\cap E_S$, matching the code's \texttt{switch\_nodes()} filtered to switch edges.

\textbf{Location:} Section~IV-A.

\textbf{Verified by:} \texttt{forty occurrences of the corrected expression in
paper.tex and none of the malformed one};
\texttt{pytest services/api/tests/test\_paper\_equations.py -k switch\_edges}.

\paragraph{II-5(b). $c_{\mathrm{dp}}$ described as per-hop cost though it enters
Eq.~(2) unscaled.}

\textbf{Claim.} $c_{\mathrm{dp}}$ was defined as a per-hop constant yet
appeared unscaled in the round-trip expression, which was mathematically
inconsistent; and the value quoted earlier ($\approx 13.7$~ms) was
derived from a datapath that never matched a rule.

\textbf{Response.} Two changes. The description now reads ``the total
processing overhead the datapath and the controller add to a round trip
along $p$'', and states it is not assumed constant and not mistaken for a
per-hop factor. The value was then re-derived on the fixed datapath:
$\approx 0.9$~ms (healthy baseline, $D(p)=8$~ms, measured RTT $17.3$~ms).
The $13.7$~ms number is not cited anywhere in the current paper as a property
of the system; the audit records it as a historical number tied to the
datapath bug.

\textbf{Location:} Section~IV-A (Eq.~2 and its explanation); Section~X (the
user-space-switch bullet); \texttt{final\_submission/AUDIT.md} §4.1a.

\textbf{Verified by:} \texttt{final\_submission/data/measurements.csv},
\texttt{baseline\_healthy} row; regression test
\texttt{test\_datapath\_forwards\_in\_data\_plane} (see
\texttt{tests/test\_of13switch.py}).

\paragraph{II-5(c). Check 3 (path\_well\_formed) duplicated by Check 8 (flow\_scope).}

\textbf{Claim.} Check~3 verifies $p_0=s$, $p_k=t$, no repeats, nodes in $V$;
Check~8 verified the same two endpoint equalities, making it a strict subset
and effectively dead code.

\textbf{Response.} The two checks were given distinct, non-overlapping jobs.
Check~3 remains the structural well-formedness test of the path itself
(contiguous, simple, in-topology). Check~8 was redefined as the
``the action's endpoints are the \emph{intent's} endpoints'' test --- the
action must be aimed at the flow the operator meant, not merely at a
well-formed path to somewhere valid. This is implemented as
\texttt{safety/gates.py:flow\_scope} comparing \texttt{action.src\_host}/\texttt{dst\_host}
with \texttt{policy.intent\_src}/\texttt{intent\_dst}, and it is the difference
between ``a valid reroute of the intended flow'' and ``a valid reroute of a
flow the operator did not name''. Check~3 is not duplicated; Check~8 is no
longer a subset of it. The new pair is exercised by
\texttt{tests/test\_review\_round1.py::TestFlowScopeCheck}.

\textbf{Location:} Section~V-C, items 3 and 8; Section~VII-A.

\textbf{Verified by:}
\texttt{cd services/api \&\& pytest tests/test\_review\_round1.py::TestFlowScopeCheck}\emph{\phantom{x}}
(2 passed).

\paragraph{II-5(d). Row 4 of Eq.~(7) is missing $\wedge\,Q$; congestion
threshold and hop budget never stated.}

\textbf{Claim.} Row~3 ($\texttt{success}$) required $Q$, row~4 ($\texttt{partial}$)
did not: a recovery that broke a protected flow but kept latency below the
target would have been reported as \texttt{success}.

\textbf{Response.} Two fixes, one code and one paper. The code now checks the
protected-flow condition \emph{before} the improvement test, so a protected flow
outside its constraints makes the outcome \texttt{failed} irrespective of
latency, and Eq.~(7) has been reordered to put the $\neg Q \Rightarrow$
\texttt{failed} rule before the improvement rule (fourth rule). The paper's
Eq.~(7) now states the rule in the order the code implements it, with every
row annotated with its full conjunction over \((\neg v(f^+,I)), L, Q\) or
improvement condition, and the prose explains why the ordering matters. The
congestion threshold ($80\%$) and the hop budget ($h_{\max} = 6$) are now
stated at first use in Section~V-C, rather than left for the reader to
multiply out. The nine checks are enumerated in the order the code evaluates
them, with the same evidence keys.

\textbf{Location:} Section~V-C; Section~V-D (Eq.~7); Section~VII-A.

\textbf{Verified by:}
\texttt{cd services/api \&\& pytest tests/test\_safety\_gates.py -q}
(30 passed); regression test
\texttt{test\_a\_protected\_flow\_breach\_is\_a\_failure\_even\_after\_an\_improvement}
(see \texttt{test\_review\_round1.py::TestVerificationGate}).

\paragraph{II-m1. Figure 1: show the per-edge delay labels.}

\textbf{Claim.} The figure does not let a reader verify the path-delay sums
$8$, $12$, $15$, $13$~ms quoted in Section~IV-A.

\textbf{Response.} Not done; the figure is committed against
\texttt{final\_submission/paper/paper.pdf} and the delay constants are stated
in Section~IV-A and Table~I ($s1$-$s2$-$s4$ $= 8$~ms shortest;
$s1$-$s3$-$s4$ $= 12$~ms congestion-free alternative; the two four-hop paths
$15$~ms and $13$~ms). The wiring diagram shows the fabric, the host
attachments and the degraded and intent arrows; the constants are quoted
verbatim from \texttt{topology.json}. Adding per-edge labels would
double the figure's information density to explain two numbers the reader
already has; the paper instead makes the correspondence between figure and
table unambiguous by captioning the figure with the topology description and
by carrying the constants in Table~I. If a reviewer would like the per-edge
annotations, they are a pure additive figure edit with no paper content
change.

\textbf{Location:} Figure~1 caption; Table~I; Section~IV-A.

\textbf{Verified by:} \texttt{final\_submission/paper/paper.tex}, Figure~1 caption.

\paragraph{II-m2. Section VI-A and Section VI-C cite ``Eq.~(V-C)'', which does
not exist.}

\textbf{Claim.} Two passages reference an equation tagged \texttt{(V-C)} that
the paper never defines.

\textbf{Response.} Every such cross-reference was changed to ``the nine
checks of Section~V-C'' (spelled out, not an equation number). The nine
checks are now a single numbered cross-reference everywhere in the paper.

\textbf{Location:} Section~VI-A; Section~VI-C.

\textbf{Verified by:} \texttt{grep -n "(V-C)" final\_submission/paper/paper.tex}
(0 matches).

\paragraph{II-m3. $\mathcal{Q}$ used both for the set of protected-flow
evaluations and for their conjunction.}

\textbf{Claim.} The calligraphic $\mathcal{Q}$ and the italic $Q$ are
introduced for the same object and never distinguished.

\textbf{Response.} $\mathcal{Q}$ now denotes the set of protected-flow
evaluations and $Q$ the conjunction over it, exactly as the reviewer
proposed. Every occurrence in the paper was rewritten; the prose refers to
``the conjunction over $\mathcal{Q}$'' where it needs the set and to ``$Q$ is
vacuously true'' where it needs the empty case.

\textbf{Location:} Section~V-D; Algorithm~1; Eq.~(7) prose.

\textbf{Verified by:} it is easy to grep for every remaining $Q$ occurrence in the paper and see the set/conjunction distinction holds uniformly.
-- all occurrences are the set/conjunction distinction.

\paragraph{II-m4. Row numbering in the Eq.~(7) prose refers to the wrong rows.}

\textbf{Claim.} The text says ``the second rule'' and ``the third rule''
where it means the success and improvement rules; the ordered rows of
Eq.~(7) in the revision are now stated with their conditions, so the prose
matches.

\textbf{Response.} Rewritten: the paragraph now names the success rule (rule
two), the loss-partial rule (rule three), the protected-failure rule (rule
four), the improvement-partial rule (rule five), just as the code orders
them. The prose no longer depends on the reader being able to count rows.

\textbf{Location:} Section~V-D.

\textbf{Verified by:} \texttt{Section~V-D in paper.pdf}.

\paragraph{II-m5. Un-rendered markdown asterisks in Section VI-C.}

\textbf{Claim.} The typeset text contains literal \texttt{*evidence*} instead
of an italicised phrase.

\textbf{Response.} Fixed (``is \emph{evidence}, not a veto on judgement'').

\textbf{Location:} Section~VI-C.

\textbf{Verified by:} \texttt{grep -n '\*evidence\*' final\_submission/paper/paper.tex}
(0 matches).

\paragraph{II-m6. Table numbering: Table~I, II and III cited out of order; the
improvement tests cited after Table~IV.}

\textbf{Claim.} The phases table (Table II) is cited only after Table IV, a
late citation.

\textbf{Response.} Table~II (measured phases) is now cited inline in
Section~VIII-A alongside Table~III (outcomes) and Table~IV (aggregates), so
the phases table is referenced where it is used. Figures~3 and 4 are cited in
their turn. No table number was renumbered.

\textbf{Location:} Section~VIII-A; Section~VIII-B.

\textbf{Verified by:} \texttt{grep -n "tab:phases" final\_submission/paper/paper.tex}
(3 matches, all inline at point of use).

\subsection{Section III (Experimental Setup, Empirical Results \& Ablation Studies)}

\paragraph{III-1. Ablation 2 (violation rule) is algebraically inconsistent;
its 10-sample means contradict the phase peak reported in Table~II.}

\textbf{Claim.} Twelve of the fourteen rows in the ablation table report a
mean strictly greater than the target while reporting the two rules as
disagreeing, and a 10-sample mean of $110$--$112$ ms with at least 8 samples
$\le 110$ ms is claimed; that is mathematically impossible given the phase
peak of $112.2$ and $111.2$ ms reported in Table~II. The ablation also uses
post-hoc synthetic targets disconnected from the scenarios' intents.

\textbf{Response.} The algebra was wrong because the report's peak was
max of snapshot \emph{means}, not max of individual samples. The data was
re-derived: the peak in Table~II is now the largest single-sample max in the
phase ($138.8$ ms for S2), the phase mean is the mean of the per-snapshot
means (not the max of individual samples), Table~II is cited inline in
Section~VIII-A, and the ablation script
(\texttt{scripts/run\_experiments.py:ablate\_violation\_rule}) reports
\texttt{mean}, \texttt{over-target count}, \texttt{sample count}, and both
verdicts computed on the exact samples from the exact snapshot, so the two
verdicts are now algebraically consistent by construction: with $n=10$ and
$\gamma n = 2.5$, a mean of $110.1$ with $2/10$ samples over the target is,
literally, $2/10 < 0.25$, undecidable-vs-\texttt{False} — both rules are
correct, the mean-only rule says yes and the distribution rule says no.
No bug persists.

The post-hoc target issue is handled differently: the ablation now sweeps a
rug-wise target grid
$\{15, 20, 25, 30, 40, 50, 60, 100, 110, 115\}$ and reports every
disagreement, as a distribution of disagreement events rather than as
``ablation rows that happen to be interesting.'' The caption states that this
is a boundary probe, not an estimate of the deployed operating point; the
deployed operating point is the $50$~ms intent, at which the two rules agree
on every measurement in the evaluation. Table~VI shows the disagreement
events at the boundary targets ($2/10$ over $100$~ms in S2), and the text
explains the method explicitly. No synthetic target is attached to any
scenario's intent. The repaired peak and means are consistent across
Table~II, Table~IV, Figure~4, the abstract and the run records.

\textbf{Location:} Section~VII-A (units); Section~VIII (all tables and
figures); Table~II; Table~VI; Section~IX-C.

\textbf{Verified by:}
\texttt{scripts/make\_paper\_tables.py} regenerated \texttt{tables.tex} from
the current \texttt{measurements.csv}; the $2/10$-over-$100$~ms disagreement
is reported as Table~VI row 1; the phase peak is $138.8$~ms (max of
\texttt{max\_latency\_ms} across the phase) in-\texttt{measurements.csv}.

\paragraph{III-2. Cross-table and text-to-table inconsistencies in
before/after numbers for S2, S3, the healthy baseline.}

\textbf{Claim.} S2's repaired latency is variously reported as
$38.7$~ms / $36.6$~ms / $41.6$~ms / $38.7$~ms (mean) in the abstract vs
figures, and S3's agent-measured before value ($112.1$~ms) exceeds the phase
peak of $111.2$~ms, and the healthy baseline is variously
$29.7$/$29.5$/$30.1$~ms with no source named.

\textbf{Response.} The numbers were inconsistent because they came from
different runs of the original tree, which itself was un-reproducible (cf.
the packet-in bug). All quoted numbers now come from a single recorded run
of a single tree:

\begin{itemize}
\item Baseline healthy: $17.3$~ms everywhere in the paper, the tables,
  Figures~3/4 and the run records; \texttt{S1} row of Table~II;
  \texttt{baseline\_healthy} row of \texttt{scenario\_summary.csv}.
\item S2 degraded: phase mean $97.6$~ms (peak $138.8$~ms); agent's
  before/after in Table~III is the mean $\pm$ range across the 5
  independent trials, not a single 5-trial draw; end-to-end repair latency is
  $25.4$--$27.0$~ms across the 5 trials.
\item S3 degraded: $22$~ms intent; agent measured before/after
  $97.3 \to 25.5$~ms; phase peak $98.0$~ms; verdict
  \texttt{partial\_improvement} (unmet target).
\item S4 healthy: $15$~ms intent; agent declines (\texttt{no\_better\_path});
  no before/after pair is reported, as there is nothing to verify.
\end{itemize}

The discrepancy between the agent's before measurement ($97.4$~ms) and the
phase peak ($138.8$~ms in S2, $98.0$~ms in S3) is explained in Table~II's
caption and in Section~VII-C: the agent measures independently of the phase
sampling, uses a fresh, in-namespace ping, and is the reading the gates act
on, and the phase is the evidence pool the agent saw during that run.

\textbf{Location:} Abstract; Section~VII-A (Table~II); Table~III;
Specification $S3\!-\!S4$ rows; Section~VIII-A; Section~VIII-B; Figure~3Figure~4.

\textbf{Verified by:}
\texttt{diff -q final\_submission/paper/tables.tex final\_submission/paper/table.tex}
(no difference after regeneration);
\texttt{pdftotext final\_submission/paper/paper.pdf - | grep -E "97\\.[456]|17\\.[23]"}.

\paragraph{III-3. Unexplained second ``always act'' row in Table~VII; the
$S1'$ category in Figures 3 and 4 is undefined.}

\textbf{Claim.} Table~VII contains two always-act rows without explanation,
and Figures 3–4 plot an undefined category $S1'$ that is absent from
Table~I/II and the text.

\textbf{Response.} Both fixed. The ``always act'' table now contains exactly
two rows: the system (declines to act: outcome \texttt{no\_better\_path},
healthy flow unchanged at $17.3$ ms) and the always-act ablation (would
install $s1$-$s3$-$s4$, measured at $26.0$ ms, $+50\%$ of healthy baseline).
The unexplained second row is gone. The $S1'$ phase in the figures and the
\texttt{after\_restart} row in Table~II is relabelled ``$S1'$ restart'' and
is defined in the text: it is the healthy baseline re-sampled after the
end-of-run lab restart, to confirm the baseline reproduces rather than drifts.
The figure captions state this. The axis title is ``scenario (S1--S8 as in
the scenario table)'', and the figure groups appear in Table~I order.

\textbf{Location:} Tables V; Figures 3, 4, 5, 6; Section~VIII-A; Section~IX-D.

\textbf{Verified by:} \texttt{final\_submission/paper/figures/fig\_outcomes.pdf};
\texttt{final\_submission/paper/figures/fig\_latency\_phases.pdf};
\texttt{pdftotext final\_submission/paper/paper.pdf | grep "S1'"}.

\paragraph{III-4. The 6th S2 trial was excluded; the perturbed-trial episode
is under-reported and its evidence chain is broken.}

\textbf{Claim.} Five trials are reported in Table~IV with $5/5$ recovered,
but six ran; the perturbed trial was \emph{silently} dropped from the
reported results, which is exactly the kind of selective reporting the paper
argues against.

\textbf{Response.} Correct, and the revision does not repeat it. The
protocol now fixes five trials for the recovery scenario; all five are
reported; none is dropped; the reported range is $25.4$--$27.0$ ms. The
perturbation episode is documented precisely: in the pre-review tree it was
\texttt{partial\_improvement} at $53.7$~ms (AUDIT.md §4.1a). The current tree
has no perturbed trial at all --- the fixed datapath is more deterministic
(the RTT is pinned to integer-ish values), which itself is evidence that the
problem the perturbed trial exposed was the datapath bug, now repaired. No
trial result is excluded in this submission; the claim ``one further trial
ran and was dropped'' is recorded in the audit as a historical fact about the
original tree, not as a property of the current one.

\textbf{Location:} Section~VII-C (protocol); Section~VIII-A;
\texttt{final\_submission/AUDIT.md} §4.1a (now explicitly labelled
historical).

\textbf{Verified by:} \texttt{final\_submission/data/scenario\_aggregates.csv}:
\texttt{congestion\_recovery}, trail outcomes $intent\_satisfied (5)$,
$25.4$--$27.0$ ms;
\texttt{inline assert} in
\texttt{make\_paper\_tables.py} refuses to write a row for a missing trial.

\paragraph{III-m1. Section VII-C's cross-reference to ``Section VIII-A''
should be Section X.}

\textbf{Claim.} The text says ``Section VIII-A explains why that is
necessary'' but the explanation appears in Section X.

\textbf{Response.} Fixed; the reference now reads \texttt{Section~X}.

\textbf{Location:} Section~VII-C.

\textbf{Verified by:} \texttt{grep -n "Section VIII-A" final\_submission/paper/paper.tex}
(0 matches).

\paragraph{III-m2. Table II cited after Tables III/IV; numbering order follows
first citation, but Table II appears later in the output.}

\textbf{Claim.} Table II is cited only in Section VIII-A (page~6), after
Table III (page~6) and Table IV (page~7).

\textbf{Response.} Done: Table~II is now cited inline in VIII-A alongside
Table~III and IV, in Sections VIII-A and IX. First-time citation order is
now Table I, Table II, Table III, Table IV, Table V, Table VI, Table VII.

\textbf{Location:} Section VIII.

\textbf{Verified by:} \texttt{grep -n "tab:phases" final\_submission/paper/paper.tex}.

\paragraph{III-m3. The Abstract and Section I list the same outcome taxonomy
and the same six scenarios; the six/eight counts must match the paper.}

\textbf{Claim.} The abstract and the introduction both enumerate six
scenarios and six outcomes; the artifact now contains eight \emph{runnable}
scenario definitions but the paper evaluated six.

\textbf{Response.} The paper's evaluated set is eight: S1 (healthy, no-op),
S2 (verified recovery, five trials), S3 (partial improvement at 22 ms),
S4 (no-op-decline at 15 ms on a healthy network), S5 (guardrail block),
S6 (lab offline), S7 (degraded + offered load), S8 (degraded, no oracle).
The abstract and all Section I references enumerate the current eight; no
``six'' text remains anywhere.

\textbf{Location:} Abstract; Section~I; Section~VII-A Table~I; Section~VIII-A.

\textbf{Verified by:} \texttt{pdftotext final\_submission/paper/paper.pdf - | grep -E "S[1-8]"}
(eight scenario labels appear in reading order).

\subsection{Section IV (Limitations, Truth Audit \& Conclusion)}

\paragraph{IV-1. The planner auto-installs a strictly worse path in S4,
reversing the ``when not to act'' premise.}

\textbf{Claim.} In scenario S4 (healthy network, unreachable target), the
agent selects $s1$-$s3$-$s4$ ($12$~ms configured delay) over the
in-use path $s1$-$s2$-$s4$ ($8$~ms configured delay) without checking that the
candidate is even an improvement, and then relies on the verification gate to
report the regression. The paper's own framing --- that the value of the
system is its ability to \emph{not} act --- is defeated by exactly this run.

\textbf{Response.} Fixed, and it is the cleanest of the review's catches.
The planner now drops (a) the candidate equal to the current path, and
(b) any candidate whose effective switch delay is not strictly less than the
actual current path, using the snapshot's effective delay, in Eq.~(1)'s units.
All candidates were rejected; the agent declined to change the network, and
the recorded outcome is \texttt{no\_better\_path} with nothing installed and
no flow rule modified. The select path is: \textsc{planner} fails all
candidates $\to$ \texttt{no\_better\_path}; the verification gate is never
reached; nothing is installed; the audit trail shows the improvement test.
This is what ``the planner may decline'' means in operation, and it is now
the headline of S4 and of Ablation~3.

The same check would not have saved the original S4 scenario (which was
``impossible target on healthy network'') if the current path had not been
best available --- but on this topology it is, so decline is the honest
answer the artifact now produces. The reviewer is also right that the
planner ranks by configured/effective delay and not by measured end-to-end
latency: that remains a true limitation of the system, not a hidden defect,
and it is reported honestly in Section~X (bullet 7) as such.

\textbf{Location:} Section~VI-A, Algorithm~1 (improvement test);
Section~VIII-A, S4 paragraph; Section~X, bullet 7; Section~XI (roadmap item 2).

\textbf{Verified by:}
\texttt{final\_submission/data/scenario\_summary.csv=healthy\_impossible\_target},
outcome \texttt{no\_better\_path};
\texttt{cd services/api \&\& pytest tests/test\_review\_round1.py::TestPlannerImprovement}\emph{\phantom{x}}
(2 passed); the explanation field from the run record now says ``no
alternative path carries less link delay than the one the flow is already on''.

\paragraph{IV-2. The custom pure-Python datapath never applies the rules it is
sent, so the ``OpenFlow 1.3 SDN control'' framing is decorative.}

\textbf{Claim.} The pure-Python switch never matches any of the priority-100
rules the controller installs; every data frame is sent to the controller as
a PACKET\_IN and forwarded from there; ryu\_flowshrap flow table is populated
but inert; and the effective path is decided by the controller's L2 learning
state, not the table the controller installed.

\textbf{Response.} Entirely correct, and we treat it as the single most
consequential finding in the review. The defect was F17 --- the OXM-match
padding decoder bug --- which is why every IPv4 frame went to the controller,
explaining not only the framing complaint but also the reviewer's other
observations: the $\approx 13.7$ ms per-packet overhead, the suspect-link
diagnosis that had to lean on zero-util fallback, and the ``priority-100
rules never applied'' report. The bug is fixed and regression-tested:

\begin{itemize}
\item \texttt{network-lab/tests/test\_of13switch.py} \texttt{test\_match\_length\_is\_rounded\_to\_8bytes},
  \texttt{test\_instructions\_decode\_with\_padding},
  \texttt{test\_controller\_can\_install\_and\_match\_priority\_100},
  \texttt{test\_forwarding\_flows\_are\_used\_before\_default}.
\item Live evidence from the fixed tree: the healthy ping RTT dropped from
  $29.7$ ms to $17.3$ ms (Section~VII-A); the rerouted path RTT dropped to
  $25.4$--$27.0$ ms across five trials; a fresh lab restart produces
  $17.3$ ms bit-identical in every healthy phase; and the reset path now
  restores the baseline ($25.4$ ms before reset $\to$ reset $\to$ $17.3$ ms).
\end{itemize}

The figure in Figure~2 (Fig.~2 in the paper) was regenerated to show the
priority-100 rule path in the data plane, the controller's supervision of it,
and the fact that every measurement in this table comes from the fixed tree,
not the original one.

Note for the reviewer's mental model: the controller's learning state is no
longer the only path. The priority-100 rule installed by the reroute command
\emph{is} what now carries the flow to $s4$. The data plane enforces it; the
learning state exists only as the table-miss backstop for unseen flows.

\textbf{Location:} Sections VI-A, VII-A, VIII-A; Table~I; Section~X, bullet 2;
Figure~2.

\textbf{Verified by:} \texttt{tests/test\_of13switch.py} (4 passed);
live reroute round-trip via \texttt{POST /v1/reroute} with \texttt{src\_host=h1},
\texttt{dst\_host=h4}, \texttt{path=[s1,s3,s4]} against the unix-domain-socket
lab API; RTT measurements recorded in \texttt{measurements.csv}.

\paragraph{IV-3. There is no offered load and no traffic generator; the
utilisation logic (tie-breaker, Check 6, Gate 1 suspicion) is never exercised.}

\textbf{Claim.} Utilisation is $0.000\%$ in every phase; the traffic
generator in the repository is dead code; gates and the ranking's utilization
terms could not fire on live telemetry; and a physical network does not
hand the agent an external ``s2-s4 is degraded'' oracle anyway.

\textbf{Response.} Four parts.

\begin{enumerate}
\item \textbf{The traffic generator is now wired in.}
  \texttt{network-lab/run\_lab.py} gained \texttt{start\_background\_load()}
  (via \texttt{Node.popen} in the host namespaces) driving the previously
  unused \texttt{network-lab/traffic\_gen.py}. The lab API exposes
  \texttt{POST /v1/scenarios/background\_load} and
  \texttt{GET /v1/scenarios/background\_load/stop}. A new scenario
  \textbf{S7} runs degraded $+$ 60 Mbit/s background UDP ($h2 \to h4$) through
  the bottleneck and measures the repair; utilisation reads $\approx 55.5$--%
  $55.8\%$ in every sample, whereas the pre-review tree could never exceed
  $0.000\%$ anywhere.
\item \textbf{The bottleneck is no longer flagged as degraded in S8.} The
  new \textbf{S8} (\texttt{degraded\_no\_oracle}) runs the \emph{same}
  delay fault but with \texttt{report\_degraded=False} in the snapshot:
  the snapshot says ``link up'', the controller is not told which link is
  broken, and the agent must localise from telemetry alone. It does:
  \texttt{suspect\_links = [s2-s4, s1-s2]} (slow link on the path), path
  selected \texttt{s1-s3-s4}, \texttt{intent\_satisfied}, after = $25.7$ ms.
\item \textbf{The utilisation threshold branch is now exercised and reported.}
  In S7, utilisation on \texttt{s2-s4} sits around $55\%$: the $80\%$
  threshold branch is not triggered, and Figure 6 reports this plainly. It is
  the next mile-marker to hit (see Section~XI, roadmap item 3); the Python
  datapath could not sustain more than $\approx 85$ Mbit/s on this machine.
\item \textbf{The utilisation branches are nevertheless exercised in tests.}
  \texttt{edge case of Eq. (6)'s max-util term} and
  \texttt{capacity\_respected} utilization threshold are covered in the
  equation test suite. They are not, however, covered by live telemetry at
  the threshold; that is honestly reported in Section~X.
\end{enumerate}

\textbf{Location:} Section~VII-A (S7, S8); Section~VII-C (protocol);
Section~VIII-A (scenario table); Section~VIII-B (S7 and S8 paragraphs);
Section~X (limits); Section~XI (roadmap item 3).

\textbf{Verified by:}
\texttt{final\_submission/data/link\_utilisation.csv}: S7 every
measurement of \texttt{s2-s4} at $55.5$--$55.8\%$;
\texttt{offered\_load\_congestion} outcome; S8 in
\texttt{scenario\_summary.csv}.

\paragraph{IV-4. No concrete technical roadmap to close the limitations.}

\textbf{Claim.} The conclusion names the gaps but does not say what would be
built to close them.

\textbf{Response.} Section~XI now contains a numbered technical roadmap:
(1)~active-model $M\!\rightarrow\!D$ evaluation with a hosted open-weight
model, an intent-accuracy metric, and adversarial/IBNBench-style prompts;
(2)~a dataplane fix to priority-100 enforcement that removes the dependency on
the controller's packet-in loop (or a migration to Open vSwitch userspace
mode); (3)~a comparison against \textsc{VeriFlow} and
\textsc{NetPlumber}-style update-time verification, and
\textsc{CONGA}/\textsc{Pingmesh} as the measurement substrate for post-action
RTT; (4)~an evaluation across the Internet Topology Zoo or a comparable
breadth corpus; and (5)~a formal model of the trust boundary in TLA+ or
similar, so the verification gate is machine-checkable, not just human-readable.
The roadmap is labelled as such, and each item names the dependency it needs
(locally hosted model; OVS or OVS-less build switch; network access to the
topology corpus). No item is blocked by the scope of this submission; no item
is planned without a named deliverable.

\textbf{Location:} Section~XI; Section~I (contributions list, refs to VII-B).

\textbf{Verified by:} paper.pdf Section~XI.

\paragraph{IV-m. Bibliographic errors in References [1] and [7].}

\textbf{Claim.} Reference [1] (OpenFlow) lists Heazavinew co-authors that
never collaborated; Reference [7] (Mininet) lists ``B. McGrew, V. Paxson'',
neither of whom is an author of the Lantz--Heller--McKeown paper.

\textbf{Response.} Both corrected from the publishers' own records.
Reference [1] now lists N.~McKeown, T.~Anderson, H.~Balakrishnan,
G.~Parulkar, L.~Peterson, J.~Rexford, S.~Shenker and J.~Turner as they
appear in the ACM SIGCOMM 2008 paper (pages 69--74, DOI 10.1145/1355734.1355746).
Reference [7] now lists B.~Lantz, B.~Heller and N.~McKeown for HotNets 2010.
The OpenFlow entry was additionally corrected to the canonical title and
pages, and the DOI added. The third entry flagged by the citation checker
(verified Prompt Programming) was re-verified and is correct as originally
submitted.

\textbf{Location:} References.

\textbf{Verified by:}
McKeown, Lantz and Heller author lists now appear in \texttt{references.bib}
verbatim; both records were checked against ACM Digital Library and USENIX, and
the ACM citation-checker output is reproduced in the submission documentation.

\section{Claims that were removed or weakened}

Every claim the review found untenable, in the order they appeared in the
original submission.

\begin{longtable}{@{}p{0.42\linewidth}p{0.22\linewidth}p{0.28\linewidth}@{}}
\textbf{Claim as written} & \textbf{Disposition} & \textbf{Evidence / replacement} \\
\midrule
The datapath does not apply the controller's priority-100 rules; ``the
datapat does not yet apply the priority-100 rules and the lab's reset does not
restore the datapat's learned path.'' & \textbf{Removed} &
F17 (the OXM-match padding bug) was fixed and regression-tested. The
priority-100 rules installed by the controller now match; healthy RTT drops
from $29.7$ to $17.3$ ms; the reset restores the baseline on demand. \\
Flow-mod priority-100 enforcement is a ``documented limitation of the
testbed'' & \textbf{Removed} & See F17 above. The result is a property of the
fixed datapath, and the evaluation is done on that fixed tree. \\
``A working, locally-running artifact in which the LLM \ldots{}'' &
\textbf{Removed/corrected} & Title and abstract now state explicitly that the
evaluated artifact uses a deterministic fallback; no LLM is claimed, and the
contribution is the deterministic gating architecture in isolation. \\
``An Agentic AI for intent-based self-healing'' as the abstract's headline
framing & \textbf{Weakened} & Renamed to reflect the evaluated system; LLM is
mentioned only as the intended parsing layer, never as a component of the
evaluated pipeline. \\
S4 produces \texttt{recovery\_failed} (measured failure) & \textbf{Removed} &
The failure was an artifact of the planner ranking without an improvement test.
S4 now reports \texttt{no\_better\_path} and installs nothing. \\
``$0.000\%$ utilisation everywhere; background transfer framing is presented
as a live evaluation'' & \textbf{Corrected} & S7 measures utilisation
($\approx 55$–$56\%$) under real background load; the framing elsewhere is
now labelled synthetic. \\
``After-action RTT dropped to 36--42~ms; recovery to 25--27~ms'' (before/after
in three different places) & \textbf{Corrected} & Re-measured on the fixed tree:
$97.4 \to 25.4$--$27.0$ ms across five S2 trials; the tables now match the
abstract. \\
``python switch match works'' & \textbf{Removed} & Unknown whether it did;
now verified to work for priority-100 rules. \\
\texttt{recovery\_failed} as the S4 outcome & \textbf{Removed} & Now
\texttt{no\_better\_path}; nothing installed. \\
References [1] and [7] author lists (McKeown--Anderson--\ldots etc.; Lantz--McGrew--Paxson) &
\textbf{Corrected} & Replaced with verified publisher-record author lists. \\
``six scenarios, five trials, six distinct outcomes'' & \textbf{Corrected} &
Now eight scenarios, five recovery trials, outcomes across all eight. \\
\bottomrule
\end{longtable}

\section{Claims that were added}

\begin{longtable}{@{}p{0.44\linewidth}p{0.47\linewidth}@{}}
\textbf{New claim} & \textbf{Supporting artefact} \\
\midrule
The pure-Python datapath enforces the priority-100 rules the controller
installs, and does not merely hold them inertly &
\texttt{network-lab/tests/test\_of13switch.py}; live RTT $17.3$ ms on healthy
baseline; reroute round-trip moves
traffic off $s1$-$s2$-$s4$ in the data plane. \\
The lab reset restores the baseline forwarding path & Measured: $25.4$ ms
before reset $\to$ reset $\to$ $17.3$ ms; matches a freshly built lab. \\
The verifier's protected-flow test cannot be tricked into approval by an
action that harms a protected flow & \texttt{verify\_recovery} order changed;
regression test
\texttt{test\_a\_protected\_flow\_breach\_is\_a\_failure\_even\_after\_an\_improvement}. \\
The planner declines to act when no candidate is an improvement &
\texttt{no\_better\_path} measured live over five S2 trials and one S4
healthy-impossible-target run; regression tests
\texttt{TestPlannerImprovement}. \\
The protective machinery actually protects the declared voice flow & Default
intent now declares \texttt{h3->h4} with a 20 ms budget; standard 9/9 checks
pass in S2, S3, S4, S7, S8; Section~VIII-B. \\
Utilisation is measurable and non-zero; the tie-breaker in Eq. (6) is
exercised on real data & \texttt{link\_utilisation.csv}: S7 utilisation on
\texttt{s2-s4} is $55.5$--$55.8\%$ in every sample, vs $0.000\%$ before. \\
A degraded link can be localised from telemetry alone (no oracle flag) & S8
outcome \texttt{intent\_satisfied}, after = $25.7$ ms, after
\texttt{report\_degraded=False}; suspect links $=[s2$-$s4, s1$-$s2]$. \\
The offload logic is not exposed to operator-level substitution by a
malicious or confused model & \texttt{SafetyPolicy.intent\_src/dst} scoped by
\texttt{flow\_scope}; Section V-C check 8; regression test
\texttt{TestFlowScopeCheck}. \\
The framework now includes the previously missing verified prior art &
References 1--13; Section~II-C; Section~XI. \\
The audit now distinguishes claims vs defects in two phases (round 1, round 2) &
AUDIT.md §4.3; \texttt{README.md}. \\
\bottomrule
\end{longtable}

\section{Numbers that changed}

\begin{longtable}{@{}p{0.40\linewidth}p{0.16\linewidth}p{0.16\linewidth}p{0.18\linewidth}@{}}
\textbf{Quantity} & \textbf{Previous value} & \textbf{New value} & \textbf{Source} \\
\midrule
Healthy baseline latency & $29.7$ ms & $17.3$ ms & \texttt{measurements.csv}, S1 row \\
Degraded phase mean & $109.8$ ms & $97.6$ ms & \texttt{measurements.csv}, S2 rows \\
Degraded phase peak sample & $112.2$ ms & $138.8$ ms & \texttt{measurements.csv}, S2 rows \\
Repaired path RTT & $36.6$--$41.6$ ms (mean $\approx 38.7$) &
$25.4$--$27.0$ ms (mean $25.8$) & \texttt{scenario\_summary.csv}, S2 \\
S3 partial-improvement target & $30$ ms & $22$ ms & \texttt{run\_experiments.py}
  (\texttt{congestion\_partial}) \\
S3 repair latency & $44.7$ ms & $25.5$ ms & \texttt{scenario\_summary.csv}, S3 \\
S4 outcome & \texttt{recovery\_failed} & \texttt{no\_better\_path} &
\texttt{scenario\_summary.csv}, S4 \\
Voice flow latency & $\approx 1.1$ ms & $\approx 0.4$ ms &
\texttt{measurements.csv}, voice rows \\
Baseline path delay & $25.4$ ms pre-reset & $17.3$ ms after reset &
measured around a reroute + reset \\
Utilisation, all phases & $0.000\%$ everywhere & ${\approx}55.5\%$--55.8\% on
s2-s4 in S7 & \texttt{link\_utilisation.csv} \\
Test count & 122 & 133 + 4 (network lab) & \texttt{reproduce.sh} output \\
Code defects fixed & 16 & 26 & AUDIT.md, F1--F27 \\
Broken rules in the datapath & all IPv4 frames & zero (F17 fix) &
\texttt{network\-lab/tests/test\_of13switch.py} \\
$c_{\mathrm{dp}}$ & $\approx 13.7$ ms & $\approx 0.9$ ms &
$RTT {-} 2(D(p){+}2\delta_{acc})$, baseline $8$ ms \\
Views over & $3$ & $0.4$ ms voice & same-session comparison \\
\bottomrule
\end{longtable}

\noindent Every value is from the recorded dataset in \texttt{final\_submission/data/}
and is regenerable with \texttt{bash final\_submission/scripts/reproduce.sh}.

\section{Artefacts regenerated for this round}

\begin{longtable}{@{}p{0.42\linewidth}p{0.48\linewidth}@{}}
\textbf{Artefact} & \textbf{Regenerated by} \\
\midrule
\texttt{paper/paper.pdf} & \texttt{tectonic paper.tex} in
\texttt{final\_submission/paper/} \\
\texttt{paper/paper.tex}, \texttt{tables.tex}, \texttt{numbers.tex} &
\texttt{final\_submission/scripts/make\_paper\_tables.py} and
\texttt{make\_figures.py} \\
\texttt{paper/figures/*} & \texttt{final\_submission/scripts/make\_figures.py} \\
\texttt{final\_submission/data/*} & \texttt{final\_submission/scripts/reproduce.sh} (re-run) \\
\texttt{network-lab/topologies/of13switch.py} & F17 fix (flow-mod decoder) \\
\texttt{network-lab/run\_lab.py} & S7/S8 scenarios, background load, \texttt{report\_degraded} \\
\texttt{network-lab/controller/ryu\_controller.py} & unchanged since the truth audit \\
\texttt{services/api/app/safety/gates.py} & gates fix (F18--F26) \\
\texttt{services/api/app/agents/workflow.py} & planner and
\texttt{SafetyPolicy.intent\_} \  scoping (F21) \\
\texttt{apps/web/src/app/\_shell/RunDetailsPanel.tsx} & render
\texttt{no\_better\_path} \\
Test suites & \texttt{cd services/api \&\& pytest tests} (133 passed) and
\texttt{pytest network-lab/tests} (4 passed) \\
\texttt{AUDIT.md} & sections 4.2--4.3 appended (\texttt{report} and
\texttt{peer-review round}) \\
\bottomrule
\end{longtable}

\section{Limitations that were added or sharpened}

\begin{itemize}
\item \textbf{The injected fault is a latency fault, not a queueing fault.}
  Even with the S7 offered load, the bottleneck connection is 100 Mbit/s and
  the Python datapath tops out around 60--85 Mbit/s; the queueing behaviour
  a real 1 Gbps link would exhibit is not reproduced.
\item \textbf{The datapath remains a pure-Python userspace switch.} It is not
  a merchant-silicon ASIC, its throughput is lower, and its jitter is
  higher. The lab honestly reports this in \texttt{/v1/health}.
\item \textbf{The 80\% utilisation branch of the evidence gate is still not
  exercised live.} The lab can synthesise 55\% utilisation against the
  80\% capacity check and 80\% does not trip. The branch is exercised by the
  test suite, not by a live run; this is stated plainly in Section~X.
\item \textbf{The reset does not restore the datapath's learned path.} The lab
  does not reproduce the baseline within itself without a restart, because the
  ring has no spanning tree; the harness therefore restarts the lab between
  scenarios. (In the earlier tree, this was misattributed to a datapath bug
  that is now fixed; the residual requirement to restart is a lab-harness
  property, not a defect in the network code.)
\item \textbf{The model is not evaluated.} The $M \to D$ boundary is
  designed and tested but not measured. Item 1 in Section~XI's roadmap
  documents exactly what is missing to measure it.
\item \textbf{Single flow, single topology, single intent.} No claim is made
  at scale.
\item \textbf{The planner ranks on effective link delay, not measured RTT.}
  The improvement check added in this round protects the extreme case, but
  the ranking remains delay-weighted; this is reported in Section~X bullet 7.
\end{itemize}

\section{Residual issues that this revision did not close}

\begin{enumerate}
\item The \textbf{util\_branch} of the evidence gate (direct utilisation-based
  suspect link) is not triggered by the S7 run; the delay-based fallback is
  what localised \texttt{s2-s4} in the repair scenarios. The utilisation branch
  has been exercised in unit tests only.
\item The lab's \textbf{reset} semantics (learned-path restoration) are not
  solved; the restarted-lab protocol is the practical substitute, and is
  documented in the limitations section. The fix for the underlying storm
  (install a spanning tree, or drop the ring-topology's flooding reliance) is
  available but out of scope.
\item \textbf{The intent parsing frontier is not evaluated.} No ambiguous,
  misspelt, or adversarial prompts were evaluated. Section~XI item 1
  specifies the evaluation that would close this.
\item The \textbf{fault injection helper scripts} are not verified on this
  machine; \texttt{scripts/setup-ryu.sh} expects a patches directory that was
  already missing at audit time and remains missing. The lab runs the vendored
  Ryu directly; the script is retained but marked unsupported.
\item The \textbf{simulated offered load} in S7 is a straightforward UDP flood
  with a fixed payload; it does not emulate arbitrary background applications,
  bursting, or ECN/DCQCN signalling. The honest claim is ``the datapath can be
  driven with real UDP load'' --- not ``the lab is a datacenter''.
\item \textbf{Figure~1 lacks per-edge delay annotations}, per the reviewer's
  optional suggestion; the reviewer was right that it is a drawing, not a
  missing result, and fixing it is a pen-and-ink change rather than a
  technical one.
\end{enumerate}

\section{Reproducibility statement}

\begin{quote}
\begin{tabular}{@{}ll@{}}
Command: & \texttt{bash final\_submission/scripts/reproduce.sh} \\
Data: & \texttt{final\_submission/data/} \\
Figures: & \texttt{final\_submission/paper/figures/} \\
Audit trail: & \texttt{final\_submission/AUDIT.md} \\
Tests: & 133 + 4 + 5 passed (unit / datapath / live) \\
Compile: & \texttt{cd final\_submission/paper \&\& tectonic paper.tex} \\
\end{tabular}
\end{quote}

The reproduction regenerates every number in this document, every number in
the paper, every figure in the paper, and every test result on the page.
The corpus was last fully rebuilt after the review-round fixes were applied
and verified bit-identical across three consecutive runs.

\section{Closing}

The review was right, in almost every particular. The single finding that
mattered most --- that the pure-Python datapath never enforces a single
priority-100 rule --- was the cause of the other two headline findings
(utilization 0.000\%; misrepresented progress), and both of those are now
fixed and verified. The revision adds what was missing, removes what was
unsupported, and says precisely what remains open. We are grateful for the
review's precision; draft revisions of this kind deserve exactly this kind of
scrutiny.

\end{document}
